\worksheettitle
  {Домашняя практика}
  {\modulemark{M01-H} Точка, прямая, отрезок и луч}

\studentnameblock

\begin{practicebox}[1. Заполните таблицу]
  \begin{center}
  \renewcommand{\arraystretch}{1.4}
  \begin{tabularx}{\textwidth}{l|c|X}
    \toprule
    Объект & Число концов & Как продолжается? \\
    \midrule
    прямая & \answerblank[18mm] & \answerblank[52mm] \\
    отрезок & \answerblank[18mm] & \answerblank[52mm] \\
    луч & \answerblank[18mm] & \answerblank[52mm] \\
    \bottomrule
  \end{tabularx}
  \end{center}
\end{practicebox}

\begin{practicebox}[2. Названия]
  \begin{enumerate}[label=\alph*)]
    \item Луч начинается в $M$ и проходит через $N$: \answerblank[25mm].
    \item Отрезок имеет концы $C,D$: два допустимых названия
      \answerblank[25mm] и \answerblank[25mm].
    \item Прямая проходит через $P,Q$: \answerblank[25mm].
  \end{enumerate}
\end{practicebox}

\begin{practicebox}[3. Построение]
  Отметьте точки $A,B,C$, не лежащие на одной прямой. Проведите прямую
  $AB$, отрезок $AC$ и луч $CB$.
  \workarea{42mm}
\end{practicebox}

\begin{warningbox}[4. Исправьте ошибку]
  «Луч $PK$ начинается в $K$». Назовите причину и запишите исправление.
  \writinglines{2}
\end{warningbox}
